\documentclass{elegantbook}
\usepackage[utf8]{inputenc}
\usepackage[T1]{fontenc}
\usepackage{geometry}
\usepackage{graphicx}
\usepackage{float}
\usepackage{hyperref}
\usepackage[english,indonesian]{babel}

\title{Optimasi Pemrosesan Template LaTeX Menggunakan Pendekatan Berbasis AI}
\author{Academic Support Agent}
\date{\today}

\begin{document}
\maketitle

\begin{abstract}

\end{abstract}

\chapter{PENDAHULUAN}


\chapter{TINJAUAN PUSTAKA}


% Original content (commented out):
% \documentclass[11pt, singlecolumn, citestyle=authoryear]{elegantbook}
\definecolor{customcolor}{RGB}{153,204,255}
\colorlet{coverlinecolor}{customcolor}

\begin{document}
\mainmatter


\chapter*{Title of the science letter goes here}


\begin{tcolorbox}[enhanced,attach boxed title to top center={yshift=-3mm,yshifttext=-1mm},colback=blue!5!white,colframe=blue!75!black,colbacktitle=blue!80!black,title=SUMMARY,fonttitle=\bfseries, boxed title style={size=small,colframe=blue!75!black} ]
  Brief summary of the science letter goes here. Science letters are informal public documents that will explore the scientific potential of next-generation (XG) gravitational-wave observatory networks. Each science letter should contain a maximum of 2 pages of text and figures, plus authors and relevant references.  They should use clear and straightforward language that is understandable for scientists from a variety of fields.

\end{tcolorbox}


\subsection*{Key question(s) and scientific context in brief} 
Explain the key science question or questions and provide a brief summary of the current status of the field on these questions.

\subsection*{Potential scientific impact of XG detectors on the key questions}
An XG network may include a combination of Cosmic Explorer, Einstein Telescope, Nemo, and upgraded detectors in current and planned ~4km-scale facilities (LIGO-US, LIGO-India, Virgo, KAGRA). Science Letters should focus on desired observational capabilities rather than specific sites or configurations. Please refer to \cite{2021arXiv210909882E} for the Cosmic Explorer Horizon Study. 
 
\subsection*{Benchmarks for XG detectors to enable the scientific impact}
Science letters are encouraged to suggest benchmarks that will ensure that the XG network is capable of the breakthrough science that they describe.  A benchmark should represent a desired observational capability, such as detection rate or sensitivity, that can be used to differentiate between potential XG networks. Science Letters should explain how reaching certain benchmarks would ensure that gravitational-wave astronomy contributes to new knowledge pathways. References should support evaluating the maturity of the benchmark values and enable a more detailed understanding where needed.

Example benchmarks could include:
\begin{itemize}
\item Characteristic strain sensitivity between 10 and 100 Hz.
\item Number of sources with 10 $\deg^2$ localization available 5 minutes before merger.
\item Number of BBH mergers detected per year with SNR$>600$. 
\item Number of BBH and BNS detected at redshift larger than 2.
\item Horizon distance for BNS observations relative to peak star formation redshift.
\item Integrated 1-year sensitivity to stochastic backgrounds with a fractional energy density in gravitational waves GW(f).
\item Sensitive range for a burst that releases EGW of  $10^{-9}$\,M\,c$^2$ at 1\,kHz.
\item Ellipticity of known pulsars that would give measurable quadrupole emission for.
\item Sensitive range for a long-duration source of some characteristic timescale and $h_{RSS}$.
\end{itemize}

\begin{tcolorbox}[enhanced,attach boxed title to top center={yshift=-3mm,yshifttext=-1mm},colback=cyan!5!white,colframe=cyan!75!black,colbacktitle=cyan!80!black,title=SCIENTIFIC IMPACT OF XG DETECTORS,fonttitle=\bfseries, boxed title style={size=small,colframe=cyan!75!black} ]
  List of foreseen science breakthroughs goes here.
\end{tcolorbox}

\subsection*{Dependencies on other multi-messenger capabilities}
Only if needed, provide a brief description of other observational capabilities that are required to address the key question(s).

\begin{tcolorbox}[enhanced,frame style image=blueshade.png,
  opacityback=0.75,opacitybacktitle=0.25,
  colback=blue!5!white,colframe=blue!75!black,
  title={\bf \centering XG DETECTOR AND NETWORK REQUIREMENTS}]
  Each science letter should outline the capabilities and timelines needed for an XG network to address a specific science question. 
\end{tcolorbox}

% template to insert figure
%
%\begin{figure}[htbp]
%  \centering
%  \includegraphics[width=0.6\textwidth]
%{figurefilenamehere}
%  \caption{caption}
%\end{figure}

\newpage
\chapter*{Authors}
\begin{center}
Author1FirstAndLastName, affiliations, email\\Author2FirstAndLastName, affiliations, email\\
Author3FirstAndLastName, affiliations, email\\
\end{center}

\printbibliography[heading=bibintoc, title=\ebibname]


\end{document}



\end{document}